% ============================================================
%  Experiment Setup — NeurIPS 2026 style
% ============================================================

\section{Experiment Setup}
\label{sec:setup}

\subsection{Task: Structured Neural Architecture Search}
\label{sec:task}

We instantiate the AutoResearch-RL pipeline on CIFAR-10 image classification.
At each research step an LLM agent proposes a neural architecture and
training-recipe configuration; a deterministic harness compiles the
configuration into a PyTorch model, trains it for a fixed number of gradient
steps, and returns the validation accuracy as the reward signal.
Unlike the original AutoResearch-RL formulation, which asks the LLM to
generate free-form Python code diffs, we restrict proposals to a
\textbf{structured JSON schema} (\texttt{ArchSpec}).
This choice (i) makes the reward fully deterministic---any valid JSON is
immediately executable---and (ii) eliminates confounds from syntax errors or
unrunnable code, allowing the comparison to isolate the contribution of search
strategy alone.

\subsection{Action Space}
\label{sec:action_space}

\texttt{ArchSpec} is a Pydantic-validated JSON object with discrete fields.
A network is specified as an ordered sequence of 1--4 convolutional stages
followed by global average pooling and a linear classifier.

\begin{table}[h]
\centering
\caption{Action space fields and valid values.}
\label{tab:action_space}
\small
\begin{tabular}{llr}
\toprule
\textbf{Field} & \textbf{Valid values} & \textbf{\# options} \\
\midrule
\multicolumn{3}{l}{\textit{Per stage (1--4 stages)}} \\
\quad \texttt{out\_channels} & 32, 64, 128, 256 & 4 \\
\quad \texttt{num\_blocks}   & 1, 2, 3           & 3 \\
\quad \texttt{norm}          & none, batch, group & 3 \\
\quad \texttt{activation}    & relu, gelu, silu   & 3 \\
\quad \texttt{downsample}    & stride, maxpool, avgpool & 3 \\
\midrule
\multicolumn{3}{l}{\textit{Network-level}} \\
\quad \texttt{use\_residual} & true, false       & 2 \\
\quad \texttt{dropout}       & 0.0, 0.1, 0.2, 0.3, 0.5 & 5 \\
\quad \texttt{scheduler}     & none, cosine, step & 3 \\
\quad \texttt{augment}       & none, basic, medium, strong & 4 \\
\midrule
\multicolumn{3}{l}{\textit{Optimizer}} \\
\quad \texttt{type}          & sgd, adam, adamw  & 3 \\
\quad \texttt{lr}            & 0.001, 0.003, 0.01 & 3 \\
\quad \texttt{weight\_decay} & 0.0, 0.001, 0.01, 0.05 & 4 \\
\bottomrule
\end{tabular}
\end{table}

The total space is approximately $6.4 \times 10^{13}$ configurations
(4-stage architectures alone: $324^4 \approx 1.1 \times 10^{10}$; combined
with $5{,}760$ non-architecture options).
With 60 research steps and $G{=}4$ candidates per step, each full run evaluates
240 configurations, covering roughly $3.7 \times 10^{-10}\%$ of the space.

\subsection{Evaluation Oracle}
\label{sec:oracle}

Each configuration is evaluated by \texttt{evaluate\_spec}, which trains the
compiled model for exactly \texttt{max\_train\_steps}$\,{=}\,300$ gradient
steps on CIFAR-10 (50k training images, 10k validation images,
batch size 128) and returns the final validation accuracy.
Fixing gradient steps rather than epochs ensures every architecture receives
an identical compute budget regardless of depth or width, making accuracy
directly comparable across configurations.
All evaluations use the same data split and a condition-specific random seed
to reduce evaluation noise.

\subsection{Experimental Conditions}
\label{sec:conditions}

We design five conditions as a sequential ablation, each isolating a single
variable relative to the previous condition.

\begin{table}[h]
\centering
\caption{Experimental conditions and their ablation targets.}
\label{tab:conditions}
\small
\begin{tabular}{clcccl}
\toprule
\textbf{ID} & \textbf{Name} & \textbf{$G$} & \textbf{History} &
\textbf{RL update} & \textbf{Ablation question} \\
\midrule
C1 & Random Search    & 1 & \texttimes & \texttimes &
    Baseline: pure random exploration \\
C2 & LLM, no history  & 1 & \texttimes & \texttimes &
    Does the LLM prior beat random? \\
C3 & LLM + history    & 1 & \checkmark & \texttimes &
    Does in-context history help? \\
C4 & LLM + history, $G{=}4$ & 4 & \checkmark & \texttimes &
    Does best-of-$G$ sampling help? \\
C5 & GRPO, $G{=}4$   & 4 & \checkmark & \checkmark (LoRA) &
    Does online RL weight update help? \\
\bottomrule
\end{tabular}
\end{table}

\paragraph{History buffer.}
Conditions C3--C5 maintain a sliding-window buffer of the $k{=}5$ most recent
experiments and the top-3 experiments by accuracy.
At each step, the full buffer is serialised into the LLM prompt so the model
can observe which configurations succeeded and which failed.

\paragraph{Parallel evaluation.}
For conditions with $G{>}1$, the $G$ candidates within a step are evaluated
in parallel using \texttt{multiprocessing} (spawn context) to avoid CUDA
context conflicts.
C5 forces sequential evaluation because simultaneously loading $G{=}4$
CNN training runs and the Qwen3-1.7B model would exceed GPU memory.

\paragraph{LLM model.}
Conditions C2--C5 use Qwen3-1.7B~\citep{qwen3} with thinking tokens enabled.
Generation uses a temperature of 0.8.
For C5, LoRA adapters ($r{=}16$, $\alpha{=}32$) are applied to the
\texttt{q\_proj} and \texttt{v\_proj} matrices of the LLM and updated
online via GRPO after each step.

\subsection{Reward Design}
\label{sec:reward}

The shaped reward passed to the GRPO update combines an accuracy signal
with an exploration bonus:
\begin{equation}
r_i =
\begin{cases}
-\delta & \text{if output is invalid JSON} \\[4pt]
f(a_i) + \lambda \cdot d_{\min}(s_i, \mathcal{H}) & \text{otherwise}
\end{cases}
\label{eq:reward}
\end{equation}
where $a_i \in [0,1]$ is the validation accuracy of candidate $i$,
$d_{\min}$ is the minimum normalised Hamming distance from $s_i$ to any
configuration in the recent history $\mathcal{H}$ (novelty bonus),
$\lambda{=}0.02$, and $\delta{=}0.1$ is the invalid-output penalty.

We experiment with two choices for the accuracy component $f(a_i)$:

\begin{itemize}
    \item \textbf{Relative:} $f(a_i) = a_i - a^*$, where $a^*$ is the
    current best accuracy. This rewards only improvement over the best
    configuration found so far.
    \item \textbf{Mixed:} $f(a_i) = \alpha \cdot a_i +
    (1-\alpha)(a_i - a^*)$, with $\alpha{=}0.5$.
    This retains an absolute accuracy signal, preventing the group reward
    from becoming uniformly negative once a strong baseline is established.
\end{itemize}

The distinction between these two formulations turns out to be consequential
for GRPO training (Section~\ref{sec:results}).
Under the relative formulation, once $a^*$ reaches a high value the rewards
for all $G$ candidates in a group are simultaneously large and negative;
the group-relative advantage
$A_i = (r_i - \bar{r}) / (\sigma_r + \varepsilon)$
has near-zero variance, effectively killing the gradient signal.
The mixed formulation avoids this degeneracy.
